\begin{solapartat}
Per puntuar aquest exercici cal tenir present que es demana una representació qualitativa, per tant, no es demana precisió sinó que quan l'estudiant dibuixi les línies de camp i equipotencials, tingui cura dels detalls que s'indiquen a continuació:

\punts{0,25} Les dues carregues estan situades correctament segons les coordenades donades a l'enunciat.

\punts{0,25} Les línies de camp apunten cap a les càrregues i totes es dirigeixen cap a una de les dues càrregues puntuals.

\punts{0,25} El nombre de línies de camp que es dirigeixen a la càrrega $q_B$ és aproximadament el doble de línies que convergeixen a $q_A$.

\punts{0,25} Les línies equipotencials són aproximadament perpendiculars a les línies de camp.

\punts{0,25} Les línies equipotencials són més properes entre elles quan ens acostem a les càrregues puntuals.

\figura[0.75]{sol_fig1.png}
\end{solapartat}

\begin{solapartat}
\punts{0,5} Com que la força que fan les dues càrregues és repulsiva, cal situar la càrrega $q_C$ entre $q_A$ i $q_B$.

\figura[0.45]{sol_fig2.png}

Com que les dues forces tenen la mateixa direcció i sentits oposats, la condició d'equilibri correspon a igualar els mòduls de les dues forces:
\[ k\,\frac{q_A q_C}{x^2} = k\,\frac{q_B q_C}{(1-x)^2}. \]
Simplificant:
\[ \frac{q_A}{x^2} = \frac{q_B}{(1-x)^2} \Rightarrow \frac{1}{x^2} = \frac{2}{(1-x)^2} \Rightarrow (1-x)^2 = 2x^2 \Rightarrow x^2 + 2x - 1 = 0 \]
L'equació anterior té dues solucions, però només la positiva és correcta:
\[ x = -1 + \sqrt{2} = 0,414\ \mathrm{m} \]

\punts{0,5} El treball que fa el camp elèctric és igual a menys la variació d'energia potencial:
\[ W = -\Delta U = -q_C\,(V_f - V_i) \]
El potencial inicial és:
\[ V_i = k\left(\frac{q_A}{x} + \frac{q_B}{1-x}\right) = 8,99\times 10^{9}\left(\frac{-2\times 10^{-6}}{0,414} + \frac{-4\times 10^{-6}}{0,586}\right) = -1,048\times 10^{5}\ \mathrm{V} \]
El potencial final és:
\[ V_f = k\left(\frac{q_A}{x} + \frac{q_B}{1-x}\right) = 8,99\times 10^{9}\left(\frac{-2\times 10^{-6}}{0,5} + \frac{-4\times 10^{-6}}{0,5}\right) = -1,079\times 10^{5}\ \mathrm{V} \]
Així el treball és:
\[ W = -\Delta U = -q_C\,(V_f - V_i) = -2,47\times 10^{-2}\ \mathrm{J} \]

\punts{0,25} El treball que fa el camp elèctric és negatiu, és a dir, el camp elèctric s'oposa al desplaçament, per tant, ha calgut aplicar una força externa per traslladar la càrrega $q_C$ des del punt inicial fins al punt final.
\end{solapartat}
